\setcounterpageref{ex}{0}
\TaoPhieuKiemTra{13}

\setcounter{ex}{0}
\setcounter{bt}{0}

\textbf{\large I. PHẦN TRẮC NGHIỆM (3,0 điểm)}

{\itshape Khoanh tròn vào chữ cái đặt trước câu trả lời đúng. Mỗi câu đúng được 0,5 điểm.}

\begin{ex} % Câu 1
Số liền sau của số bé nhất có hai chữ số là
\choice
{$10$}
{\True $11$}
{$12$}
{$9$}
\loigiai{
Số bé nhất có hai chữ số là $10$, số liền sau là $11$.\\
\textbf{Đáp án đúng: B}
}
\end{ex}

\begin{ex} % Câu 2
Trong phép cộng $18 + 70 = 88$, số hạng thứ nhất là
\choice
{\True $18$}
{$70$}
{$88$}
{$8$}
\loigiai{
Số hạng thứ nhất là $18$.\\
\textbf{Đáp án đúng: A}
}
\end{ex}

\begin{ex} % Câu 3
Bao bột mì nặng $25\text{ kg}$, người ta đã dùng hết $5\text{ kg}$. Trong bao còn lại số bột mì là
\choice
{$30\text{ kg}$}
{$15\text{ kg}$}
{$10\text{ kg}$}
{\True $20\text{ kg}$}
\loigiai{
Còn lại: $25 - 5 = 20\text{ kg}$.\\
\textbf{Đáp án đúng: D}
}
\end{ex}

\begin{ex} % Câu 4
Đường gấp khúc gồm $3$ đoạn thẳng dài $2\text{ cm}, 3\text{ cm}, 4\text{ cm}$. Độ dài đường gấp khúc là
\choice
{$8\text{ cm}$}
{$10\text{ cm}$}
{\True $9\text{ cm}$}
{$7\text{ cm}$}
\loigiai{
$2 + 3 + 4 = 9\text{ cm}$.\\
\textbf{Đáp án đúng: C}
}
\end{ex}

\begin{ex} % Câu 5
So sánh: $8 + 6 \dots 7 + 7$.
\choice
{$>$}
{\True $=$}
{$<$}
{Không so sánh được}
\loigiai{
$8 + 6 = 14$; $7 + 7 = 14$. Vì $14 = 14$ nên điền dấu $=$.\\
\textbf{Đáp án đúng: B}
}
\end{ex}

\begin{ex} % Câu 6
Tính: $9 + 4 - 8 = \dots$
\choice
{$4$}
{$6$}
{$7$}
{\True $5$}
\loigiai{
$9 + 4 = 13$; $13 - 8 = 5$.\\
\textbf{Đáp án đúng: D}
}
\end{ex}

\vspace{4mm}
\textbf{\large II. PHẦN TỰ LUẬN (7,0 điểm)}

{\itshape Thí sinh trình bày chi tiết bài làm vào giấy kiểm tra.}

\begin{bt}[1,0 điểm]
Tính nhẩm:
\begin{multicols}{2}
\begin{enumerate}[a)]
  \item $9 + 7$
  \item $7 + 8$
  \item $15 - 6$
  \item $17 - 7$
\end{enumerate}
\end{multicols}
\loigiai{
\begin{itemize}
  \item $9 + 7 = 16$
  \item $7 + 8 = 15$
  \item $15 - 6 = 9$
  \item $17 - 7 = 10$
\end{itemize}
}
\end{bt}

\begin{bt}[1,5 điểm]
Đặt tính rồi tính:
\begin{multicols}{3}
\begin{enumerate}[a)]
  \item $49 + 46$
  \item $81 - 40$
  \item $42 + 50$
\end{enumerate}
\end{multicols}
\loigiai{
\begin{itemize}
  \item $49 + 46 = 95$
  \item $81 - 40 = 41$
  \item $42 + 50 = 92$
\end{itemize}
}
\end{bt}

\begin{bt}[1,0 điểm]
Tính giá trị biểu thức kèm đơn vị đo:
\begin{enumerate}[a)]
  \item $29\text{ kg} + 8\text{ kg} - 10\text{ kg}$
  \item $38\text{ }l - 14\text{ }l + 7\text{ }l$
\end{enumerate}
\loigiai{
\begin{enumerate}[a)]
  \item $29\text{ kg} + 8\text{ kg} - 10\text{ kg} = 37\text{ kg} - 10\text{ kg} = 27\text{ kg}$.
  \item $38\text{ }l - 14\text{ }l + 7\text{ }l = 24\text{ }l + 7\text{ }l = 31\text{ }l$.
\end{enumerate}
}
\end{bt}

\begin{bt}[1,5 điểm]
Bạn An có $28$ cây bút chì. Bạn Bình có nhiều hơn bạn An $5$ cây bút chì. Hỏi bạn Bình có bao nhiêu cây bút chì?
\loigiai{
Số cây bút chì bạn Bình có là:
\[ 28 + 5 = 33 \text{ (cây bút chì)} \]
\textbf{Đáp số:} $33$ cây bút chì.
}
\end{bt}

\begin{bt}[1,0 điểm]
Một thùng chứa đầy $48\text{ }l$ nước mắm. Buổi sáng cửa hàng đã bán ra $17\text{ }l$. Hỏi trong thùng còn lại bao nhiêu lít nước mắm?
\loigiai{
Trong thùng còn lại số lít nước mắm là:
\[ 48 - 17 = 31\text{ }l \]
\textbf{Đáp số:} $31\text{ }l$ nước mắm.
}
\end{bt}

\begin{bt}[1,0 điểm]
Tìm số có hai chữ số giống nhau mà tổng hai chữ số của nó bằng $14$.
\loigiai{
Vì hai chữ số giống nhau và có tổng bằng $14$ nên mỗi chữ số là $7$ (do $7 + 7 = 14$).\\
Vậy số đó là $77$.\\
\textbf{Đáp số:} $77$.
}
\end{bt}
